\section{FLOP Estimate of Parcae}
\label{sec:flop-estimate}


In standard, fixed-depth architectures, a common means to approximate the number of FLOPs used in training is $C = 6ND$ from \citet{kaplan2020scalinglawsneurallanguage}, where $N$ is the number of parameters and $D$ is the number of tokens used in training. However, looped architectures differ from traditional models in that they exhibit the notion of \emph{effective parameters} $\hat N$ (e.g., for a model that is a single layer with $N$ parameters, if it is looped ten times, then it has an effective parameterization of $\hat N = 10 N$). Furthermore, as Parcae uses truncated backpropagation through depth, the effective parameters can thus be decoupled into two types: $\hat N_1$, which are effective parameters that \emph{\textbf{are not backpropagated}} through, and $\hat N_2$, which are effective parameters that \emph{\textbf{are backpropagated}} through. Thus, following \citet{kaplan2020scalinglawsneurallanguage}, we can formulate the effective FLOPs of Parcae as $C = (2\hat N_1 + 6 \hat N_2)D$, which further matches the setup of \citet{mcleish_retrofitted_recurrence}. Like \citet{mcleish_retrofitted_recurrence}, we exclude embedding parameters from $\hat N$, however, we do include unembedding parameters in $\hat N$ similar to \citet{nanochat}. Lastly, we additionally include an estimate for attention FLOPs following \citet{chowdhery2022palmscalinglanguagemodeling,nanochat}.
